\documentclass[11pt]{article}
\usepackage[a4paper,margin=1in]{geometry}
\usepackage{amsmath,amssymb,amsthm,booktabs}
\usepackage[T1]{fontenc}
\usepackage{lmodern}
\usepackage[hidelinks]{hyperref}
\setlength{\emergencystretch}{3em}

\theoremstyle{plain}
\newtheorem{theorem}{Theorem}
\theoremstyle{remark}
\newtheorem{remark}[theorem]{Remark}

\title{A Counterexample to Conjecture 00000000013}
\author{%
  Feiyu\thanks{Independent researcher.}
  \and
  Claude (Opus 5)\thanks{Anthropic.}%
}
\date{2026-09-12}

\begin{document}
\maketitle

\begin{abstract}
Conjecture 00000000013 asserts that for every $n\ge3$ the interval
$[2^n,2^{n+1})$ contains a prime with exactly three binary ones. It fails at
$n=8$: all twenty-eight such integers are composite. The refutation is
formalized in Lean~4 against Mathlib, with the conjecture itself stated in the
formalization.
\end{abstract}

\section{Statement}

Conjecture 00000000013 claims that for every integer $n\ge3$ there is a prime
$p$ with
\[
  2^n\le p<2^{n+1}
\]
whose binary expansion contains exactly three $1$'s. We disprove it by
exhibiting $n=8$.

\section{The counterexample}

\begin{theorem}
No prime $p$ with $2^8\le p<2^9$ has exactly three $1$'s in its binary
expansion.
\end{theorem}

\begin{proof}
An integer of $[2^8,2^9)$ has its top bit at position $8$, so it has exactly
three binary ones precisely when it is of the form
\[
  2^8+2^a+2^b,\qquad 0\le b<a<8 .
\]
There are $\binom82=28$ such integers. Each is composite, as witnessed by the
proper divisor listed beside it:

\begin{center}
\begin{tabular}{rrrrrrrr}
\toprule
$N$ & div & $N$ & div & $N$ & div & $N$ & div\\
\midrule
259 & 7  & 261 & 3  & 262 & 2  & 265 & 5\\
266 & 2  & 268 & 2  & 273 & 3  & 274 & 2\\
276 & 2  & 280 & 2  & 289 & 17 & 290 & 2\\
292 & 2  & 296 & 2  & 304 & 2  & 321 & 3\\
322 & 2  & 324 & 2  & 328 & 2  & 336 & 2\\
352 & 2  & 385 & 5  & 386 & 2  & 388 & 2\\
392 & 2  & 400 & 2  & 416 & 2  & 448 & 2\\
\bottomrule
\end{tabular}
\end{center}

Every entry in the second column is a divisor $d$ with $1<d<N$, so each $N$ is
composite. Twenty-one of the twenty-eight are even and exceed $2$; the seven
odd ones are
\[
  259=7\cdot37,\quad 261=3^2\cdot29,\quad 265=5\cdot53,\quad
  273=3\cdot7\cdot13,
\]
\[
  289=17^2,\quad 321=3\cdot107,\quad 385=5\cdot7\cdot11 .
\]
\end{proof}

Since the conjecture asserts the existence of such a prime for \emph{every}
$n\ge3$, the single value $n=8$ refutes it.

\begin{remark}
$n=8$ is not isolated. A direct search finds that $n=25$ and $n=32$ also admit
no such prime: of the $\binom{25}2=300$ and $\binom{32}2=496$ candidates in the
respective intervals, none is prime. These were checked by computation and are
not part of the Lean development, where a kernel-level decision procedure over
intervals of that size is not practical.
\end{remark}

\section{Formal verification}

The accompanying Lean~4 project formalizes the refutation against Mathlib.

\begin{itemize}
  \item \texttt{ConjectureHolds} --- the conjecture itself:
    \[
      \forall n\ge3,\ \exists p,\
      2^n\le p\ \wedge\ p<2^{n+1}\ \wedge\
      \texttt{(Nat.bits }p\texttt{).count true}=3\ \wedge\ \texttt{Nat.Prime }p .
    \]
    The digit condition counts the \texttt{true} entries of the binary
    expansion \texttt{Nat.bits}\,$p$ supplied by Mathlib, rather than a
    parametrization assumed to describe it, and primality is Mathlib's
    \texttt{Nat.Prime}.
  \item \texttt{card\_candidates} --- the interval contains exactly $28$
    integers with three binary ones, matching $\binom82$.
  \item \texttt{no\_three\_one\_prime\_at\_eight} --- no such integer is prime,
    decided by exhausting all $256$ integers of $[2^8,2^9)$.
  \item \texttt{conjecture\_00000000013\_false} --- $\lnot$\,\texttt{ConjectureHolds},
    by instantiating the conjecture at $n=8$.
\end{itemize}

Because the conjecture is stated in the formalization, the final theorem is a
literal formalization of its negation rather than of facts that merely entail
it.

The project builds with \texttt{lake build} on
\texttt{leanprover/lean4:v4.33.1} with Mathlib \texttt{v4.33.1}. It contains no
\texttt{sorry}, no \texttt{admit} and no \texttt{native\_decide}, and
introduces no axioms of its own:
\texttt{\#print axioms conjecture\_00000000013\_false} reports exactly the
three standard axioms \texttt{propext}, \texttt{Classical.choice} and
\texttt{Quot.sound}.

\end{document}
